\documentclass[letterpaper,11pt]{article}
\usepackage[spanish,mexico]{babel}
\usepackage[utf8]{inputenc}
\usepackage{graphicx}
\usepackage{array}
\usepackage{float}
\usepackage{xcolor}
\usepackage{listings}
\usepackage{tocloft}
\renewcommand{\cftsecleader}{\cftdotfill{\cftdotsep}}
\usepackage[
    left=25mm,
    right=25mm,
    top=35mm,
    bottom=30mm,
    headheight=77pt,
    papersize={21.59cm,27.94cm}
]{geometry}
\usepackage{setspace}
\onehalfspacing
\usepackage{parskip}
\usepackage{microtype}
\usepackage{fancyhdr}
\pagestyle{fancy}
\fancyhf{}
\fancyhead[L]{\includegraphics[height=1.5cm]{portada_img/der.png}}
\fancyhead[R]{\includegraphics[height=1.5cm]{portada_img/izq.png}}
\fancyfoot[C]{\thepage}
\renewcommand{\headrulewidth}{0.4pt}
\renewcommand{\footrulewidth}{0pt}
\usepackage{hyperref}
\hypersetup{
    colorlinks=true,
    linkcolor=blue,
    citecolor=blue,
    filecolor=magenta,
    urlcolor=cyan
}
\lstdefinestyle{output}{
    backgroundcolor=\color{black!5},
    basicstyle=\ttfamily\footnotesize,
    numbers=none,
    frame=single,
    rulecolor=\color{black!30},
    breaklines=true,
    showstringspaces=false,
    xleftmargin=0.5cm,
    xrightmargin=0.5cm
}
\usepackage{xurl}
\newcommand{\captura}[2]{%
\begin{figure}[H]\centering
\IfFileExists{image/#1}{%
\includegraphics[width=\textwidth,height=0.57\textheight,keepaspectratio]{image/#1}}{%
\fbox{\parbox[c][2cm][c]{0.9\textwidth}{\centering Evidencia pendiente de la prueba en Linux:\\ \texttt{\detokenize{#1}}}}}
\caption{#2}\end{figure}}

\begin{document}
\hypersetup{pageanchor=false}
\thispagestyle{empty} % Sin encabezado ni pie en la portada

% --- ENCABEZADO DE PORTADA (TABLA DE 3 COLUMNAS) ---
\begin{center}
    % La tabla usa 'm' para centrar verticalmente las imágenes y el texto
    \begin{tabular}{m{0.15\textwidth} m{0.60\textwidth} m{0.15\textwidth}}
        % Logo Izquierdo (UNAM)
        \includegraphics[width=2cm]{portada_img/der.png} & 
        
        % Título Central
        \centering\Huge\textbf{Universidad Nacional Autónoma de México} &
        
        % Logo Derecho (Facultad)
        \raggedleft\includegraphics[width=2cm]{portada_img/izq.png}
    \end{tabular}
\end{center}

\vspace{0.3cm}

% --- INFORMACIÓN DEL REPORTE ---
\begin{center}
    {\LARGE \textbf{Facultad de Ingeniería}}\\[1cm]

    {\LARGE \textbf{Alumno:}}\\[0.3cm]
    {\Large Lara Hernandez Angel Husiel - 320060829} \\[1cm]


    {\LARGE \textbf{Arquitectura Cliente-Servidor}} \\[1cm]

	{\LARGE Grupo: 02 - Semestre: 2027-1} \\[1cm]
    
    {\LARGE \textbf{Tarea 5:}}\\[0.5cm]

    {\LARGE {Sniffer de red en Linux.}}\\[1cm]

    {\LARGE \textbf{Profesor:}}\\[0.2cm]
    {\Large MTRO. CARLOS ALBERTO ROMAN ZAMITIZ}\\[1cm]

    {\Large \textbf{Fecha de realización:}}\\[0.2cm]
    {\Large 30 de octubre del 2026}
\end{center}
\newpage
{\hypersetup{linkcolor=black}\tableofcontents}
\thispagestyle{empty}
\newpage
\setcounter{page}{1}
\hypersetup{pageanchor=true}
\singlespacing

\section{Objetivo de la práctica}
Para el caso de esta tarea se utiliza el programa \texttt{sniffer.c} y el par de programas 04 del tema 3. La intención es observar cómo el servidor envía una frase al cliente y cómo el sniffer permite leer ese mismo texto dentro del tráfico TCP del puerto 4898.

\textbf{Estado del documento:} la prueba del sniffer en Linux y las capturas están pendientes. Los resultados que se describen más adelante son los esperados; todavía no corresponden a una captura realizada.

\subsection{Funcionamiento del sniffer:}
Ahora bien, el sniffer utiliza \texttt{AF\_PACKET} y \texttt{SOCK\_RAW} para recibir las tramas con su cabecera de enlace. Primero revisa Ethernet, después IPv4 y por último TCP. Para calcular dónde comienza el mensaje toma en cuenta las longitudes de las cabeceras IP y TCP, ya que pueden variar.

El filtro acepta paquetes cuyo puerto de origen o destino sea 4898. Después imprime el payload cuando contiene datos. Lo que se busca comprobar es que la frase se puede leer porque esta comunicación no utiliza cifrado. El sniffer observa la comunicación; no sustituye al servidor ni al cliente.

\subsection{Códigos utilizados:}
Se prepararon copias de los archivos del laboratorio en \texttt{Tarea5/code}. En el sniffer se cambió el filtro de 8080 a 4898, como indica la tarea. Además de eso, se agregaron comprobaciones de longitud y memoria antes de analizar las cabeceras. No se realiza reconstrucción de flujos TCP ni de fragmentos IPv4.

En el servidor 04 se cambió el puerto fijo por un argumento de la terminal. En el cliente y el servidor se revisan los errores y se conserva la frase original: \texttt{El mensaje del servidor es: Conexion exitosa!}. Los archivos originales del laboratorio no se modificaron.

\section{Equipo para realizar la prueba}
La opción recomendada es Ubuntu dentro de UTM en la Mac, utilizando tres terminales normales de la misma máquina virtual. También se puede usar Linux nativo o una VM Ubuntu en una PC. En todos esos casos se puede apuntar el cliente a \texttt{127.0.0.1}.

Por otra parte, este sniffer no se ejecuta directamente en macOS, porque \texttt{AF\_PACKET} y \texttt{linux/if\_packet.h} pertenecen a Linux. No es necesario utilizar otro equipo ni conectarse al servidor del profesor para esta práctica.

\section{Procedimiento en Ubuntu}
Los comandos se escriben uno por uno. Primero se copia la carpeta \texttt{Tarea5} al repositorio dentro de Ubuntu. La ruta de los ejemplos siguientes se debe ajustar si el repositorio está en otra ubicación.

Si faltan el compilador o las herramientas de red, se pueden instalar con:
\begin{lstlisting}[style=output]
sudo apt update
sudo apt install build-essential iproute2 net-tools
\end{lstlisting}

\subsection{Terminal 1: activar y ejecutar el sniffer:}
Antes de realizar el cambio se consulta el estado de \texttt{lo}, para poder regresarlo a su condición inicial al terminar. Después se activa el modo promiscuo y se comprueba con \texttt{ifconfig}, como se pide en la guía.
\begin{lstlisting}[style=output]
cd ~/Documentos/Cliente-Servidor/Tarea5/code
ip link show lo
sudo ip link set lo promisc on
ifconfig lo
gcc -Wall -Wextra sniffer.c -o sniffer
sudo ./sniffer
\end{lstlisting}
El comando \texttt{ifconfig lo} debe mostrar \texttt{PROMISC}. Se deja el sniffer abierto antes de iniciar la comunicación. Los permisos de administrador permiten crear el socket en crudo.

\captura{configuracion.png}{Configuración de la interfaz lo y comprobación del modo promiscuo.}

\subsection{Terminal 2: levantar el servidor:}
En una segunda terminal se compila el servidor y se pasa el puerto 4898 como argumento:
\begin{lstlisting}[style=output]
cd ~/Documentos/Cliente-Servidor/Tarea5/code
gcc -Wall -Wextra programa04_servidor_internet.c -o programa04_servidor_internet
./programa04_servidor_internet 4898
\end{lstlisting}
El servidor queda esperando una conexión. Al conectarse el cliente, envía la frase y termina. Para repetir la prueba se debe volver a ejecutar el servidor.

\subsection{Terminal 3: ejecutar el cliente:}
Por último, se abre otra terminal y se ejecuta:
\begin{lstlisting}[style=output]
cd ~/Documentos/Cliente-Servidor/Tarea5/code
gcc -Wall -Wextra programa04_cliente_internet.c -o programa04_cliente_internet
./programa04_cliente_internet 127.0.0.1 4898
\end{lstlisting}
Al presionar Enter se realiza la conexión. En este ejemplo no se escribe un nombre desde el teclado, ya que el servidor envía automáticamente la frase original del programa 04.

\section{Resultado esperado y evidencia}
Para el caso de esta prueba, el cliente debe imprimir la frase que recibe del servidor. Al mismo tiempo, el sniffer debe mostrar la MAC, la IP de origen, los puertos y el texto dentro del payload. En el segmento enviado por el servidor, el puerto de origen debe ser 4898.

Como se utiliza loopback, las MAC pueden aparecer con ceros y el sniffer puede mostrar copias del mismo paquete. Además de eso, el mensaje podría llegar en varios segmentos. Esto se debe revisar en la captura real, sin asumir que cada impresión corresponde a un envío nuevo de la aplicación.

La guía solicita una captura de las tres terminales. Se deben acomodar juntas y conservar visibles la frase enviada, la frase recibida y el payload del sniffer. El cliente y el servidor pueden haber terminado; las ventanas se mantienen abiertas con sus salidas.

\captura{prueba_tres_terminales.png}{Evidencia pendiente: sniffer, servidor y cliente en tres terminales de Linux.}

\subsection{Finalización de la prueba:}
Después de guardar la captura se detiene el sniffer con Ctrl+C. Si \texttt{lo} inicialmente no estaba en modo promiscuo, se restaura de la siguiente manera:
\begin{lstlisting}[style=output]
sudo ip link set lo promisc off
ip link show lo
\end{lstlisting}
Si el modo promiscuo ya estaba activo antes de la práctica, se conserva ese estado.

\section{Consideraciones para varios equipos}
La alternativa más sencilla es colocar el servidor y el sniffer en el mismo Linux y ejecutar el cliente desde otro equipo de la red, utilizando la IP del servidor. En ese caso se activa el modo promiscuo en la interfaz de red correspondiente, no en \texttt{lo}.

No obstante, si el sniffer se coloca en una tercera máquina, activar el modo promiscuo no garantiza que el switch le entregue el tráfico entre los otros dos equipos. La red virtual o un puerto espejo debe permitir que esas tramas lleguen al observador. Por este motivo se recomienda completar primero la prueba en una sola VM.

\section{Comprobaciones previas}
El cliente y el servidor se compilaron en la Mac sin advertencias y se probaron en localhost. Se recibieron los 46 bytes de la frase original y ambos terminaron correctamente. Esta comprobación verifica la comunicación TCP; no reemplaza la ejecución del sniffer ni la evidencia de las tres terminales en Linux.

\section{Referencias}
\begin{enumerate}
\item Guía de la Tarea 5.\\
\url{https://cromanzamitiz.github.io/arquitectura-cliente-servidor/tema3/Tarea_5.pdf}
\item Material de teoría recomendado por el profesor.\\
\url{https://cromanzamitiz.github.io/arquitectura-cliente-servidor/tema3/Clase_Sniffer_Redes_y_Cliente_Servidor.pdf}
\item Manual de Linux: sockets de paquetes y permisos.\\
\url{https://man7.org/linux/man-pages/man7/packet.7.html}
\end{enumerate}
\end{document}
